
Neural networks exploit spurious correlations in training data. Gradient-based interventions have been proposed for single-run robustification without group labels. This work tests a foundational assumption underlying Progressive Gradient Orthogonalization (PGO): that early training gradients capture spurious feature directions. Experiments on the Waterbirds benchmark reveal a measurement limitation---accumulating one gradient per epoch produces representations too low-rank to distinguish spurious from core directions. With 10 gradient samples in a 25.6-million-parameter space, both spurious and core alignments measure approximately 0.05, indistinguishable from each other. This null result stems from insufficient accumulation density rather than absence of the hypothesized effect. The findings establish a minimum requirement for valid gradient subspace analysis: multi-batch accumulation is necessary to achieve sufficient rank for meaningful directional measurement. The PGO hypothesis remains unverified pending implementation of adequate gradient sampling.

